\section{The Newton threshold and the exact zero}\label{inf:sec-newton}

The closing argument uses a threshold for the residual and two
polynomial operator bounds in the same spaces. It requires the
actual integer-split gap already obtained.

\begin{theorem}[Polynomial Newton threshold]\label{inf:newton-threshold}
Suppose fixed exponents $b_{\mathrm{inv}},m_{\mathrm{mat}},d_{\mathrm{nl}}\ge0$ have been chosen independently
of the centre accuracy $N$. Assume, uniformly in $0<\tau=|\theta|\le20$
and the admissible splits, for all sufficiently large $p$, that
in the real Banach spaces $\Xc_p,\Cc$ of
Definition~\ref{inf:spaces} a centre $z_0$ satisfies
\begin{align*}
 D\Fnl(z_0)&=J_D\Qmat+\mathcal V_{\mathrm{res}},&
 \nrm{\Fnl(z_0)}&\le C_N\theta^2p^{2-N},\\
 \nrm{\mathcal V_{\mathrm{res}}}&\le C_N\theta^2p^{3-N},&
 \nrm{J_D^{-1}}&\le C_Np^{b_{\mathrm{inv}}}/(c_{\mathrm{gap}}\tau p^{-5/3}),\\
 \nrm{\Qmat^{-1}}&\le C_Np^{m_{\mathrm{mat}}},&
 \sup\nrm{D^2\Fnl}&\le C_Np^{d_{\mathrm{nl}}}.
\end{align*}
Here $J_D$ and $\Qmat$ are bounded isomorphisms, and the last bound
holds on the valid closed centre ball considered below, on which
$\Fnl$ is $C^2$.
Put
\begin{equation}\label{inf:threshold-exponents}
 B_{\mathrm{Newt}}=b_{\mathrm{inv}}+m_{\mathrm{mat}}+5/3,\qquad
 L_{\mathrm{Newt}}=B_{\mathrm{Newt}}+d_{\mathrm{nl}}+4.
\end{equation}
Assume that the accuracy order satisfies
\begin{equation}\label{inf:threshold-accuracy}
 N\ge\lceil2B_{\mathrm{Newt}}+d_{\mathrm{nl}}+8\rceil.
\end{equation}
For $p$ above a threshold depending on the fixed constants and
$N$, the closed ball of radius
$\mathfrak r_p=\tau p^{-L_{\mathrm{Newt}}}$ about $z_0$ contains a
real zero of $\Fnl$, provided that ball is inside the domain of the
map. The threshold is uniform for $0<\tau\le20$.
\end{theorem}
\begin{proof}
Let $\AN=\Qmat^{-1}J_D^{-1}$. This isomorphism
$\Cc\to\Xc_p$ satisfies
$\nrm{\AN}\le C_Np^{B_{\mathrm{Newt}}}/\tau$.
The three quantities controlling the contraction are
\begin{align}
 \frac{\nrm{\AN\Fnl(z_0)}}{\mathfrak r_p}
 &\le C_Np^{B_{\mathrm{Newt}}+L_{\mathrm{Newt}}+2-N}
                                      \le C_Np^{-2},
 \label{inf:newton-gate-one}\\
 \nrm{I-\AN D\Fnl(z_0)}
 &\le C_N\tau p^{B_{\mathrm{Newt}}+3-N}\le20C_Np^{-5},
 \label{inf:newton-gate-two}\\
 \sup\nrm{\AN D^2\Fnl}\,\mathfrak r_p
 &\le C_Np^{B_{\mathrm{Newt}}+d_{\mathrm{nl}}-L_{\mathrm{Newt}}}
                                      =C_Np^{-4}.
 \label{inf:newton-gate-three}
\end{align}
The exact factorization gives
$I-\AN D\Fnl(z_0)=-\AN\mathcal V_{\mathrm{res}}$, proving
the second estimate. The exponents follow from
\eqref{inf:threshold-exponents} and \eqref{inf:threshold-accuracy}.
Increase $p$ until each quantity is at most $1/4$.
For $\mathcal T(z)=z-\AN\Fnl(z)$, Taylor's integral formula
bounds its displacement from the centre on the closed ball by
$\mathfrak r_p/4+\mathfrak r_p/4+\mathfrak r_p/8<\mathfrak r_p$.
The same formula for its derivative gives Lipschitz constant at
most $1/4+1/4=1/2$.
The ball is complete, so Banach's contraction theorem supplies
a fixed point. Since $\AN$ is injective, its fixed-point equation
implies $\Fnl(z)=0$. All the smallness conditions are uniform for
$0<\tau\le20$, as the displayed powers show.
\end{proof}

\begin{corollary}[Instantiation for the centre family]
\label{inf:Newton-zero}
Fix the constants in the following order:
\begin{equation}\label{inf:constant-order}
 \begin{aligned}
 C_0&\longrightarrow C_{\mathrm{rel}}\longrightarrow
 E_{\mathrm{inv}}\longrightarrow M=M_{20E_{\mathrm{inv}}}\\
 &\longrightarrow (D_{\mathrm{sep}},K_{\mathrm{head}})
 \longrightarrow b_{\mathrm{inv}}=2M+3
 \longrightarrow N=4M+40.
 \end{aligned}
\end{equation}
Here $C_0$ and $C_{\mathrm{rel}}$ are those of
\eqref{inf:uniform-delta} and \eqref{inf:relative-two-norms}.
For every sufficiently large $p\in\mathcal P_{\mathrm{bd}}$, choose a split
from Theorem~\ref{inf:integer-gap} and the order-$N$ centre at
that split. There is an exact real zero $z_*=(g_*,f_*)$ with
\begin{equation}\label{inf:Newton-correction}
 \nrm{z_*-z_0}_{\Xc_p}\le\tau p^{-L_{\mathrm{Newt}}},\qquad
 B_{\mathrm{Newt}}=2M+32/3,\quad L_{\mathrm{Newt}}=2M+62/3.
\end{equation}
\end{corollary}
\begin{proof}
The centre gap follows from Theorem~\ref{inf:integer-gap} and
Proposition~\ref{inf:centre-gap-transfer} at the now fixed order.
The coefficient inverse has exponent $b_{\mathrm{inv}}$ by
Theorem~\ref{inf:coefficient-inverse}.
The material inverse and second derivative have exponents
$m_{\mathrm{mat}}=d_{\mathrm{nl}}=6$ by \eqref{inf:Qmat-bounds} and \eqref{inf:p6-second}.
The residuals and factorization are
\eqref{inf:residual-moments}, \eqref{inf:residual-material-bound},
and \eqref{inf:material-factorization}.
Thus all hypotheses of Theorem~\ref{inf:newton-threshold} hold
in exactly $\Xc_p=\Gc_p\times\Ec_1$ and $\Cc_{p,2}$.
The choice $N=4M+40$ exceeds
$\lceil2B_{\mathrm{Newt}}+14\rceil$.
The three gate exponents in this instance are
\[
 B_{\mathrm{Newt}}+L_{\mathrm{Newt}}+2-N=-20/3,\qquad
 B_{\mathrm{Newt}}+3-N=-2M-79/3,\qquad
 B_{\mathrm{Newt}}+6-L_{\mathrm{Newt}}=-4.
\]
The second gate retains its factor $\tau$; the first and third
have no residual negative power of it.
Since $\nrm{f_N}_{\Ec_1}=O_N(\tau p^{-2})$ and
$\mathfrak r_p\to0$ uniformly for $\tau\le20$, the closed ball
is eventually in the product region
\eqref{inf:nonlinear-working-region}: its source displacement is
at most one and its total shape norm is at most $1/4$. Increase the threshold to meet all preceding estimates
for this one fixed $N$ and apply the threshold theorem.
The spaces, centre and operators are real by
Definition~\ref{inf:spaces} and their defining formulas, so the
contraction acts on the real closed ball.
\end{proof}
